\section{不变的核心能力}

\subsection{AI 无法替代的技能}

尽管工具和工作流在快速变化，但 Mihail Eric 在 Week 1 提出的核心理念在未来十年仍然成立：

\begin{importantbox}{穿越技术周期的核心能力}
\begin{enumerate}
    \item \textbf{业务理解}：理解用户需求和商业价值——AI 不知道``该做什么''
    \item \textbf{架构思维}：系统级的设计决策需要权衡无数的 trade-off
    \item \textbf{品味（Taste）}：区分好代码和坏代码、好设计和坏设计的能力
    \item \textbf{代码阅读能力}：审查和理解 AI 生成的代码比自己写代码更重要
    \item \textbf{沟通与协调}：无论是管理人还是管理 Agent，清晰的沟通是基础
\end{enumerate}
\end{importantbox}

\subsection{给新一代开发者的建议}

基于整门课程的洞察，对未来开发者的几点建议：

\begin{enumerate}
    \item \textbf{学习 AI 工具，但不要依赖它们}：理解底层原理（Week 1--2）比熟练使用某个工具更持久
    \item \textbf{投资于阅读代码的能力}：在 AI 编码时代，审查能力比编写能力更有价值（Week 7）
    \item \textbf{练习委派和多线程工作}：这是 Agent Manager 模式的核心技能（Week 3--4）
    \item \textbf{理解安全的基本原则}：AI 引入了新的攻击面，安全素养是不可协商的（Week 6）
    \item \textbf{拥抱快速实验}：正如 Mihail 在 Week 1 所说，``没有成熟的范式，每个人都在摸索''
    \item \textbf{为 6 个月后的模型构建}：Boris Cherny 的建议（Week 4）适用于所有工具和工作流的选择
\end{enumerate}

\subsection{未来团队形态的演化}

如果把整门课的结论映射到组织层面，未来的软件团队很可能出现新的分层：

\begin{table}[H]
\centering
\begin{tabular}{p{0.2\textwidth} p{0.26\textwidth} p{0.28\textwidth}}
\toprule
团队层次 & 主要工作方式 & 竞争优势来源 \\
\midrule
个人开发者层 & 借助 AI IDE 和 Agent 完成端到端任务 & 试错速度、跨职能执行力 \\
小团队层 & 多人各自管理多个 Agent 并协同集成 & 规划能力、接口边界和 review 机制 \\
平台团队层 & 维护统一的工具链、规范、评测与安全基线 & 组织规模化复制能力 \\
\bottomrule
\end{tabular}
\caption{AI 软件工程时代的团队形态变化}
\end{table}

\begin{knowledgebox}{未来的组织优势会更多来自系统设计，而不是单点高手}
当个人编码效率被 AI 大幅放大后，真正拉开差距的往往不再是“谁写得更快”，而是“谁能把工具、流程、安全和审查体系组装成一个稳定系统”。这也是整门课最后回到组织与流程视角的原因。
\end{knowledgebox}

\subsection{本章小结}

技术变化越快，基础能力越重要。业务理解、架构思维、品味、代码阅读和沟通是穿越 AI 技术周期的永恒技能。

